\chapter{A Jangada}

\lettrine{O}{\space Buddha apontou para a forma de ver} as coisas como elas são; isto é o que
queremos dizer com ``iluminação''. Vendo a forma como as coisas realmente são,
não estamos condenados a viver numa dimensão da qual não há saída. Existe uma
saída clara dessa dimensão de miséria; uma forma muito precisa. Então o Buddha
disse: ``Eu ensino apenas duas coisas: sofrimento e o fim do sofrimento.''

O Budismo é uma religião desconcertante para os ocidentais porque não possui uma
posição doutrinária. Não faz nenhuma declaração doutrinária sobre a realidade
última ou seja o que for: há apenas sofrimento e o fim do sofrimento. Isso deve
ser realizado agora; para concretizar o fim do sofrimento temos de admitir e
entender realmente o que é o sofrimento, pois o problema não é com o sofrimento,
mas com a ilusão e a ganância. E realmente temos que entender o sofrimento --
então de acordo com o sermão das Quatro Nobres Verdades, o sofrimento deve ser
compreendido. Existe sofrimento. Deve ser compreendido.

Na nossa vida diária aqui em Amaravati, percebemos quando estamos a sofrer.
Podemos culpar o clima, as pessoas ou qualquer outra coisa, mas isso não é o que
interessa; porque mesmo que alguém esteja a tratar"-nos mal, o mundo é assim
mesmo. Às vezes as pessoas tratam"-nos bem, às vezes tratam"-nos mal por causa
desta preocupação mundana por condições; mas o sofrimento é algo que criamos.

Num mosteiro, tentamos agir de forma responsável para não causar sofrimento a
ninguém intencionalmente. Estamos aqui para nos encorajarmos uns aos outros em
direcção à responsabilidade moral, à cooperação, à bondade e compaixão. Essa é a
nossa intenção.

Às vezes perdemo"-nos -- explodimos uns com os outros ou fazemos coisas que não
são muito simpáticas -- mas essa não é a nossa intenção; estes são os momentos
descuidados. Eu comporto"-me de maneira moral não apenas para meu próprio
benefício, para a minha própria prática, mas por respeito a vós e ao Sangha,
para a comunidade que nos rodeia: ser alguém que vive dentro das restrições dos
preceitos morais.

Então a minha intenção diz respeito ao meu relacionamento convosco, na direcção
de \emph{metta}, bondade e compaixão, alegria, calma e serenidade. Pelo menos a
intenção de cada um de nós é fazer o bem e abster"-se de fazer o mal. E isso
ajuda"-nos a olhar para o sofrimento que criamos numa comunidade com este tipo de
objectivo -- pois muitos de vocês realmente sofrem aqui. E isto deve ser
compreendido. É a Primeira Nobre Verdade, \emph{dukkha}, o sofrimento de não
conseguir o que queremos; o sofrimento de as coisas não serem como queremos, da
separação do que gostamos; o sofrimento de ter que fazer o que não queremos
fazer, de termos que nos controlar quando queremos ser desenfreados.

Penso o quão fácil é eu criar uma noção dos outros na minha mente. ``As monjas
são assim, os Anagārikas são assim, os bhikkhus são assim'' e assim por diante.
Podemos ter estes preconceitos: ``As mulheres são assim, os homens são assado;
os americanos são assim, e os ingleses são assado.'' Podemos acreditar nisso,
mas essas são percepções da mente, noções que surgem e cessam. E, no entanto,
podemos criar imenso sofrimento através delas. ``Este aqui não vem aos cânticos
da manhã'' ou ``aquele não está a fazer a sua parte no trabalho'', e ``este
pensa que eles são muito importantes'' ou seja o que for, mas o ponto importante
é o sofrimento, o \emph{dukkha}, porque quando temos isso, criamos desespero nas
nossas mentes. Ficamos chateados e indignados e tudo isso nos conduz ao
desespero. Se não entendemos \emph{dukkha} aqui, então não vamos entendê"-lo onde
quer que estejamos: seja em Londres, Bangkok ou em Washington DC; no topo de uma
montanha ou num vale; com as pessoas boas ou as más pessoas. Portanto, é muito
importante observar o sofrimento para se conhecer o \emph{dukkha}.

Existem três compreensões lúcidas da Primeira Nobre Verdade: existe
\emph{dukkha}; este deve ser compreendido; este foi compreendido. É assim que a
compreensão lúcida funciona; o seu reconhecimento; é algo para se compreender;
começar a perceber quando o compreendemos. Portanto, essas são as três
compreensões lúcidas da Primeira Nobre Verdade.

A Segunda Nobre Verdade é a origem de \emph{dukkha}: existe uma origem; é por
causa do cobiçar do desejo. A segunda compreensão lúcida da Segunda Nobre
Verdade é que este apego ao desejo -- esta identificação como ``eu'' e ``meu''
com o desejo, este perseguir o desejo -- deve ser abandonado, deixando"-o como
está. Então, a terceira compreensão lúcida da Segunda Nobre Verdade é: o desejo
foi abandonado -- através da prática. Abandonou"-se \emph{dukkha}.

Há a primeira compreensão lúcida em cada uma das verdades: \emph{pariyatti} --
uma observação de que existe sofrimento, sua origem e assim por diante. Depois
existe o \emph{paṭipatti} ou a compreensão lúcida na prática. O que fazemos.
Como praticamos. E então a terceira compreensão lúcida é o \emph{paṭivedha} ou a
sabedoria. Foi compreendido; abandonou"-se.

Ora, quando ocorre aquela compreensão lúcida: ``A origem do sofrimento foi
abandonada'', existe o conhecimento desse resultado -- na verdade, o desapego.
Sabemos o que é não estarmos apegados a algo. Como este relógio. Isto é segurar
o relógio; é assim. E \emph{agora} estou ciente de como é não segurar o relógio.
Se estou a segurar coisas e não presto atenção, então nem sequer me apercebo que
não estou a segurar as coisas. Quando não há cobiça, não me apercebo dela. Uma
pessoa verdadeiramente ignorante e descuidada está tão apanhada pela avidez, que
embora não esteja o tempo inteiro na posse de algo, o hábito é de tal ordem que
só se apercebe quando se encontra no encalce de algo. Como agora, muitos de
vocês só se sentem totalmente vivos quando estão cheios de ganância ou raiva de
uma forma ou de outra. Portanto, abrir mão disso pode ser bastante assustador
para as pessoas; quando abrem a mão das coisas, já não se sentem vivas.

Há um enorme investimento em ser uma pessoa. Até mesmo a visão de ``Eu tenho mau
feitio; eu tenho muita raiva'' pode ser um tipo de presunção. Se estou zangado
sinto"-me muito vivo. O desejo sexual faz com que o ``eu'' se sinta vivo -- e é
por isso que há tanta obsessão com sexo nas vidas europeias modernas. E quando
não há desejo sexual, não há raiva, eu quero adormecer. Não sou nada. Quando não
há consciência plena de todo, apenas temos de procurar mais prazer sensual --
comer algo, beber qualquer coisa, usar drogas ou ver qualquer coisa na TV, ler
qualquer coisa ou fazer algo perigoso. Podemos infringir a lei só porque é
excitante fazê"-lo.

Agora imaginem tentar fazer as pessoas passarem um fim de semana segurando
apenas um relógio, percebendo como é segurar um relógio! Que perda de tempo, eu
podia estar aí fora a aterrorizar a polícia; podia estar numa discoteca -- com
luzes estroboscópicas, música estridente nos meus ouvidos, com cannabis e L.S.D.
e uísque! Por comparação, estar atento ao modo como as coisas são e deixar
apenas de distrair a mente soa muito doloroso.

Esta noite vamos sentar"-nos em meditação até à meia"-noite. É uma oportunidade de
observar mais plenamente como é estar sentado; como é quando a mente está cheia
de pensamentos e quando não há pensamentos; quando há sofrimento e quando não há
sofrimento. Se têm uma ideia de que ficarem sentados até à meia"-noite vai ser
sofrimento, então já se comprometeram a sofrer até à meia"-noite.

Mas se começarem a examinar essa mesma ideia ou medo ou dúvida na vossa mente
pelo que é, podem ver quando isso está presente e quando não está presente. Se
estiverem a sofrer, então não estão a pensar que há sofrimento. Então têm essa
sensação de sofrimento e estão agarrados à ideia ``Estou a sofrer e tenho que me
sentar direito e estou cansado''. Portanto, a Primeira Nobre Verdade: ``há
sofrimento; o sofrimento deve ser compreendido'' e isto através de uma admissão,
um reconhecimento e uma compreensão.

A compreensão lúcida da Segunda Nobre Verdade é abandoná"-la, deixá"-la em paz --
não tirem conclusões de uma sessão nocturna. Estas são percepções. Elas são na
verdade, nada: se usarem a situação para reflexão e contemplação de quando há
sofrimento -- então não há qualquer sofrimento. Estamos cientes de termos este
pensamento, de procurar este pensamento ou de não procurar este pensamento.
Podemos pegar ou largar coisas, sabendo como usar estas coisas em vez de termos
uma obsessão cega para agarrar ou rejeitar. Eu posso pousar o relógio, mas não
preciso de deitá"-lo fora, certo? Não é que segurar o relógio esteja errado, a
não ser que haja ignorância sobre isso. Uma pessoa está ciente do agarrar e do
não agarrar, do segurar e do não segurar.

A Terceira Nobre Verdade -- existe a cessação do sofrimento. Quando se abre mão
de algo e se percebe isso, os nossos hábitos tornam"-se os nossos professores.
Quando largamos o sofrimento, o sofrimento cessa. ``Há cessação e isso deve ser
compreendido'' -- esta é a segunda compreensão lúcida da Terceira Nobre Verdade.
E esta é a nossa prática: compreender a cessação, perceber quando o sofrimento
cessa. Não é que tudo vá desaparecer, mas o sentimento de sofrimento e ``eu
sou'' cessa. Isso não é para ser acreditado, mas para ser compreendido -- e
depois há a terceira compreensão lúcida: que a cessação foi compreendida.

Isto conduz à compreensão lúcida da Quarta Nobre Verdade a respeito do Caminho
óctuplo, o Caminho para sair do sofrimento. Estas compreensões lúcidas ligam"-se
umas às outras. Não se trata de perfazer primeiro uma e depois a outra -- elas
suportam"-se. À medida que compreendemos lucidamente o desapego, à medida que
compreendemos a cessação, logo existe o Entendimento Correcto e o resto segue"-se
a isto -- o desenvolvimento da sabedoria ou \emph{paññā}.

Agora, não vejam isto como algo que lida apenas com questões muito profundas e
importantes, porque se trata do aqui e agora, da maneira como as coisas são. Não
estamos apenas a pensar em situações extremas para trabalhar, mas apenas
sentarmo"-nos, estarmos de pé, andarmos, deitarmo"-nos, respirarmos, sentindo"-nos
seres normais, vivermos num ambiente moral da forma como as coisas são. Não
precisamos ir ao Inferno para realmente ver o sofrimento; não estamos à procura
disso.

Podemos criar o Inferno em Amaravati, não porque Amaravati seja o Inferno, mas
porque o criamos com todos os tipos de coisas miseráveis da nossa mente e este é
o sofrimento com o qual podemos trabalhar. É apenas o sofrimento do reino humano
normal, onde as nossas intenções são evitar fazer o mal, fazer o bem,
desenvolver a virtude e sermos amáveis. Ainda temos aqui sofrimento suficiente
para contemplarmos estas Quatro Nobres Verdades em seus doze aspectos.

Podemos memorizá"-las; assim, onde quer que estejamos, teremos algo para
contemplar. Eventualmente, abrimos mão de todas estas coisas porque elas também
não são um fim em si próprias, mas são como ferramentas para serem usadas.
Aprendemos a usar estas ferramentas e, quando terminámos, não precisamos de
continuar a apoiar"-nos nelas.

Significando isto, o Buddha referiu"-se ao seu ensinamento como uma jangada que
podemos fazer com as coisas ao nosso redor. Não precisamos ter uma lancha
especial, ou submarino, ou transatlântico de luxo. Uma jangada é algo que
fazemos com as coisas à volta apenas para atravessar para a outra margem. Não
estamos a tentar fazer um veículo super impressionante; podemos usar o que está
ao nosso redor para a iluminação. A jangada é para nos transportar através do
mar da ignorância e quando chegarmos à outra margem, podemos largá"-la -- o que
não significa que tenhamos que deitá"-la fora.

Esta ``outra margem'' também pode ser uma ilusão, porque ``a outra margem'' e
esta margem são realmente a mesma. É apenas uma alegoria. Nunca deixamos
realmente a outra margem, sempre estivemos na outra margem de qualquer maneira;
e a jangada é algo que usamos para nos lembrar que realmente não precisamos de
uma jangada. Portanto, não há absolutamente nada a fazer, excepto estarmos
conscientes, sentarmo"-nos, estar de pé, andar, deitar, comer a nossa comida,
respirar -- todas as oportunidades como humanos para fazer o bem. Temos esta
oportunidade maravilhosa no reino humano de sermos bons, amáveis, generosos, de
amar, servir e ajudar os outros. Esta é uma das qualidades mais bonitas do ser
humano.

Podemos decidir não fazer o mal. Não temos que matar, mentir, roubar,
distrair"-nos e drogar"-nos ou perdermo"-nos em humores e sentimentos. Podemos
livrar"-nos de tudo isso. É uma oportunidade maravilhosa de na forma humana nos
abstermos do mal e de fazer o bem -- não a fim de acumular mérito para a próxima
vida, mas porque esta é a beleza da nossa humanidade. Ser humano pode ser uma
experiência cheia de alegria em vez de uma tarefa pesada.

E então, quando contemplamos isto, começamos realmente a apreciar este
nascimento numa forma humana. Sentimo"-nos gratos por ter esta oportunidade de
viver com o nosso professor, o Buddha, e nossa prática, o Dhamma; e viver
na Sangha.

Sangha representa a comunidade humana unificada em virtuosidade e restrição
moral; é a força da alma do reino humano. Aquilo que é verdadeiramente
benevolente na humanidade tem o seu efeito nos aspectos morais que suportam o
reino humano.

Portanto, todos os seres sencientes são beneficiados por isso. Como seria se
houvesse apenas uma humanidade egoísta, sendo cada um por si, fazendo exigências
intermináveis, sem nos importarmos de todo uns com os outros? Seria um lugar
terrível para se viver. Portanto, não fazemos isso; nós permanecemos na Sangha,
uma convivência onde vivemos juntos dentro de uma convenção que incentiva a
moralidade e o respeito mútuo. Isto é para reflexão, para contemplação; temos
que saber isso por nós mesmos; ninguém o pode compreender por nós. Temos que nos
motivar e não depender de algo externo para nos impelir ou segurar.

Precisamos até mesmo de abandonar a nossa necessidade de ser inspirados. Temos
que desenvolver a força até não mais precisarmos de qualquer tipo de inspiração
ou encorajamento de ninguém -- porque inspiração não é sabedoria, é? Vocês ficam
eufóricos e -- ``Ajahn Sumedho é maravilhoso'' -- e passado um tempo já não
ficam eufóricos comigo, e então: ``Ajahn Sumedho é decepcionante; ele
decepcionou"-me.'' A inspiração é como comer chocolate: tem um gosto bom e é
muito atraente, mas não alimenta; só energiza momentaneamente e é tudo o que
pode fazer. Portanto, não é aconselhável depender de outras pessoas que vivem da
maneira que queremos e que nunca nos decepcionam.

É extremamente importante desenvolver a compreensão lúcida através da prática,
porque a inspiração esgota"-se simplesmente -- e se nós estamos agarrados e cegos
por ela, então estamos prontos para uma terrível desilusão e mágoa. Há tanto
disso com diferentes figuras de gurus carismáticas a ensinarem à volta do mundo.
Não é equilibrado, certo? Por mais intoxicados que possamos ficar com o carisma
de outra pessoa, não se consegue manter. Portanto, isso inevitavelmente envolve
cair nalgum estado inferior.

O caminho do ser consciente é, no entanto, sempre apropriado ao tempo e ao
lugar, à maneira como as coisas estão nos seus aspectos bons e maus. Então o
sofrimento não depende do mundo ser bom ou mau, mas de como estamos dispostos a
usar a sabedoria neste momento presente. A saída para o sofrimento é
\emph{agora}, ser capaz de ver as coisas como elas são.

\photoPage{image-006-aj-chah.jpg}{O Venerável Ajahn Chah no Manjusri Institute}{Foto por Thubten Yeshe}
